% !TEX program = xelatex
\documentclass[11pt,a4paper]{article}
\usepackage{../../../00_common/preamble}

\title{2022年度 (AY2022) 同志社大学大学院 生命医科学研究科\\
医工学コース 入学試験問題「数学」解答例}
\author{}
\date{}

\begin{document}
\maketitle
\vspace{-3em}

%======================================================================
\section*{第1問}

\subsection*{前提整理}

$\sinh(x^2+y)$ の Maclaurin 展開を全次数 $3$ 次まで求める.
$u=x^2+y$ とおき $\sinh u=u+\dfrac{u^3}{6}+\cdots$ を用いる（$x^2$ は $2$ 次,$y$ は $1$ 次）.

\subsection*{計算}

$\sinh u=u+\dfrac{u^3}{6}+\cdots$ に $u=x^2+y$ を代入する.
$u=x^2+y$ はそのまま($\le3$ 次),$u^3=(x^2+y)^3$ の $\le3$ 次は $y^3$ のみ
（$3x^2y^2$ 以降は $4$ 次以上）.よって
\[
\ans{\
\sinh(x^2+y)=y+x^2+\frac{y^3}{6}+(\text{4次以上})\ }
\]

検算:$x=0$ で $\sinh y=y+\dfrac{y^3}{6}+\cdots$ に一致.
$y=0$ で $\sinh(x^2)=x^2+\cdots$ の $x^2$ 係数 $1$ も一致.$\checkmark$

%======================================================================
\section*{第2問}

\subsection*{前提整理}

球 $D:\ x^2+y^2+z^2\le a^2$ 上で $\displaystyle\iiint_D y^2\,dV$ を求める.
対称性を使い,球座標で計算する.

\subsection*{計算}

\paragraph{Step 1: 対称性で簡約する.}
$x,y,z$ の対称性から $\iiint_D x^2=\iiint_D y^2=\iiint_D z^2$ なので
\[
\iiint_D y^2\,dV=\frac13\iiint_D(x^2+y^2+z^2)\,dV=\frac13\iiint_D r^2\,dV.
\]

\paragraph{Step 2: 球座標で積分する.}
$dV=r^2\sin\theta\,drd\theta d\phi$ より
\[
\iiint_D r^2\,dV
=\int_0^{2\pi}\!\!d\phi\int_0^{\pi}\!\!\sin\theta\,d\theta\int_0^{a}r^4\,dr
=2\pi\cdot2\cdot\frac{a^5}{5}=\frac{4\pi a^5}{5}.
\]
したがって
\[
\ans{\ \iiint_D y^2\,dx\,dy\,dz=\frac13\cdot\frac{4\pi a^5}{5}=\frac{4\pi a^5}{15}\ }
\]

検算:$3$ 成分の和 $\iiint(x^2+y^2+z^2)=3\cdot\frac{4\pi a^5}{15}=\frac{4\pi a^5}{5}$ が
$\iiint r^2$ と一致する.$\checkmark$

%======================================================================
\section*{第3問}

\subsection*{前提整理}

$\vv_1=(1,0,1,0)^\top,\ \vv_2=(0,1,1,0)^\top,\ \vv_3=(2,0,0,1)^\top,\ \vv_4=(1,1,1,1)^\top$
の1次独立性を判定し,独立なら Gram--Schmidt で正規直交化する.

\subsection*{計算}

\paragraph{Step 1: 1次独立の判定.}
$4$ ベクトルを列にもつ行列の行列式は
\[
\det\begin{pmatrix}1&0&2&1\\0&1&0&1\\1&1&0&1\\0&0&1&1\end{pmatrix}=-3\neq0.
\]
よって $4$ ベクトルは\textbf{1次独立}（$\R^4$ の基底）である.

\paragraph{Step 2: Gram--Schmidt.}
直交化を順に行う:
\[
\begin{aligned}
\vu_1&=\vv_1,&\|\vu_1\|^2&=2,\\
\vu_2&=\vv_2-\frac{\vv_2\cdot\vu_1}{2}\vu_1=\bigl(-\tfrac12,1,\tfrac12,0\bigr)^\top,&\|\vu_2\|^2&=\tfrac32,\\
\vu_3&=\vv_3-\vu_1+\tfrac23\vu_2=\bigl(\tfrac23,\tfrac23,-\tfrac23,1\bigr)^\top,&\|\vu_3\|^2&=\tfrac73,\\
\vu_4&=\vv_4-\vu_1-\tfrac23\vu_2-\tfrac57\vu_3=\bigl(-\tfrac17,-\tfrac17,\tfrac17,\tfrac27\bigr)^\top,&\|\vu_4\|^2&=\tfrac17.
\end{aligned}
\]
正規化して
\[
\ans{\
\bm{e}_1=\frac{(1,0,1,0)^\top}{\sqrt2},\
\bm{e}_2=\frac{(-1,2,1,0)^\top}{\sqrt6},\
\bm{e}_3=\frac{(2,2,-2,3)^\top}{\sqrt{21}},\
\bm{e}_4=\frac{(-1,-1,1,2)^\top}{\sqrt7}\ }
\]

検算:各 $\|\bm{e}_i\|=1$,かつ任意の $i\neq j$ で $\bm{e}_i\cdot\bm{e}_j=0$
（例:$\bm{e}_3\cdot\bm{e}_4=\frac{1}{\sqrt{147}}(-2-2-2+6)=0$）.正規直交系をなす.$\checkmark$

%======================================================================
\section*{第4問}

\subsection*{前提整理}

\[
A=\begin{pmatrix}a&1&0\\0&a&0\\0&0&b\end{pmatrix}
\]
の対角化可能性を判定する.固有値ごとに固有空間の次元を調べる.

\subsection*{計算}

\paragraph{Step 1: 固有値.}
上三角なので固有値は対角成分 $a$（重複度 $2$）と $b$.

\paragraph{Step 2: $\lambda=a$ の固有空間.}
\[
A-aI=\begin{pmatrix}0&1&0\\0&0&0\\0&0&b-a\end{pmatrix}.
\]
$b\neq a$ のとき $\operatorname{rank}(A-aI)=2$ なので幾何的重複度は $3-2=1$,
代数的重複度 $2$ に満たない.
$b=a$ のときも $A-aI=\begin{psmallmatrix}0&1&0\\0&0&0\\0&0&0\end{psmallmatrix}$ で
$\operatorname{rank}=1$,幾何的重複度 $2<3$.

\paragraph{Step 3: 判定.}
いずれの場合も $\lambda=a$ に対する1次独立な固有ベクトルが重複度ぶん取れない
（左上に Jordan 細胞 $\begin{psmallmatrix}a&1\\0&a\end{psmallmatrix}$ を含む）ので
\[
\ans{\ A\ \text{は対角化}\textbf{不可能}\ }
\]

検算:$a=2,b=5$ とすると $A-2I$ の核は $\operatorname{Span}\{(1,0,0)^\top\}$ の $1$ 次元のみで,
固有値 $2$ の固有ベクトルが $1$ 本しか取れず,対角化に必要な $3$ 本に満たない.$\checkmark$

\end{document}
